\textbf{Proof.}
In the homogeneous specialization, class \(1\) has its unique minimum squared loading on direction \(v_1\), and class \(2\) has its unique minimum squared loading on direction \(v_2\).  Applying \(\mathrm{prop:block-specific-single-class-reduction}\) therefore gives product laws with independent symmetric signs satisfying
\[
R_1=\frac{4c_1}{n^2}=r_1^\star,
\qquad
R_2=\frac{4c_2}{n^2}=r_2^\star.
\tag{1}
\]
The gains are positive by \(\mathrm{thm:quadruple-menu-frontier}\), so each one-class retained-gain value is \(1\).

For the two-class problem, that theorem gives
\[
\lambda^\star=
\begin{cases}
(a-2b)/(a+b),&a\ge2b,\\[4pt]
(2b-a)/\{2(a+b)\},&a\le2b.
\end{cases}
\]
In the first regime the denominator exceeds the numerator by \(3b>0\); in the second it exceeds the numerator by \(3a>0\).  The relevant numerator is nonnegative in each regime.  Hence \(0\le\lambda^\star<1\) for every \(a,b>0\).  No direction-overlap condition is used, because both oracle laws are permissible boundary points of \(\Pi_n\). \(\square\)